% !TEX root = ../reporte_final.tex
\section*{Resumen ejecutivo}
\addcontentsline{toc}{section}{Resumen ejecutivo}

Este documento reporta una investigación corta de minería de datos sobre el dataset \emph{Palmer Penguins} \citep{horst2020palmer,gorman2014ecological}. Se compara la capacidad de tres algoritmos de clustering no supervisado --- K-Means \citep{macqueen1967some}, DBSCAN \citep{ester1996density} y Agglomerative Clustering con linkage de Ward \citep{ward1963hierarchical} --- para recuperar automáticamente la estructura de especies a partir únicamente de las cuatro variables morfométricas del dataset, sin utilizar la etiqueta de especie durante el entrenamiento.

\paragraph{Aportes principales}
\begin{itemize}
  \item Flujo de minería de datos completo en un único notebook reproducible: definición del problema, obtención de datos, exploración, preprocesamiento, experimentación, evaluación y conclusiones.
  \item Comparación cuantitativa de los tres algoritmos con dos métricas complementarias: \emph{Silhouette Score} (interna) y \emph{Adjusted Rand Index} (externa, contra la etiqueta verdadera de especie) \citep{rand1971objective,hubert1985comparing,rousseeuw1987silhouettes}.
  \item Visualizaciones vectoriales: \emph{pairplot} discriminado por especie, matriz de correlación, método del codo, proyección PCA bidimensional y matriz de confusión.
  \item Práctica plena de Ciencia 2.0: datos con licencia CC-0, código bajo control de versiones, semillas fijadas (\texttt{RANDOM\_STATE=42}), \texttt{requirements.txt} y notebook ejecutable de principio a fin sin intervención manual.
\end{itemize}

\paragraph{Resultado del experimento}
Sobre 342 observaciones limpias del dataset (tras eliminar valores faltantes en las variables numéricas y en la especie), \textbf{Agglomerative Clustering con linkage de Ward} obtiene el mejor ajuste a la estructura biológica conocida, con un \emph{Adjusted Rand Index} de \textbf{0.916}, un \emph{Silhouette Score} de 0.454 y, tras mapear cada cluster a su especie mayoritaria, una \textbf{exactitud del 96.8\,\%} (331/342). K-Means queda muy próximo (ARI=0.793) y DBSCAN, pese a maximizar la métrica interna (Silhouette=0.557), fusiona las especies \emph{Adelie} y \emph{Chinstrap} en un único cluster denso y produce 33 puntos catalogados como ruido, lo que penaliza su ARI a 0.584.

\paragraph{Defensibilidad académica}
El experimento es totalmente reproducible: tanto el notebook como las semillas, los hiperparámetros y los hashes del dataset están versionados en el repositorio del proyecto. La sección~6 incluye la tabla comparativa de métricas y la sección~7 las conclusiones con limitaciones y trabajo futuro. Los apéndices recogen el código fuente esencial y la checklist de reproducibilidad.
\newpage
